\scanpage{371}
\mspage{\deleted{\unsure{109}}\unsure{110|119}}
% archive sheet number: 370
d'hypothèses supplémentaires draconiennes sur $\deleted{\unsure{C}}T$, la catégorie $C$ ``nue'' n'offre pas d'intérêt particulier, \inserted{à part son $\pi_0$ (qui est \unsure{donné|déterminé} aussi par $C_\infty$ et par $H_{\unsure{C|T}}$).} Il n'est pas dit, \deleted{\unsure{certes}}, que $\unsure{\underline{\mathrm{Tg}}|\underline{T_g}}$ ne satisfasse pas à des ``conditions draconiennes''. Mais comme $C$ ``nu'' ne \deleted{se} s'insère pas dans un formalisme sympa (p. ex. on n'a pas d'interprétation géométrique de $\mathrm{Hom}_{\mathrm{Cat\,hot}}(S,C)$, ou de $\mathrm{Hom}_{\mathrm{Hot}}(S,C)$), je crois que pour la suite je peux en faire abstraction. Exit \unsure{$\to$|a)} !

\hfill$\backslash$\hfill\mbox{}

\hfill 5.11.90

J'ai quand même passé la journée entière hier à essayer de comprendre la relation de $C$ et de son type dans $\mathrm{Hot}$, à celui de $C_\infty$. En effet, comme on a finalement une flèche dans $\mathrm{Hot}$
\[
C_\infty\longrightarrow C,
\]
déduite du diagramme dans $(\mathrm{Hot})$
\diagram{\begin{tikzcd}[ampersand replacement=\&]
C_{\mathbb N^\circ}\arrow[r,"\sim"]\arrow[d]\& C_\infty\\
C\&
\end{tikzcd}}
de prouver que $C_\infty\longrightarrow C$ est un isom. dans $\mathrm{Hot}$
